% Source: https://artofproblemsolving.com/community/c3h1617861_arithmetic_or_geometric
% Stable identity: aops-topic-1617861-q01
A sequence is arithmetic if consecutive terms have a constant difference, and geometric if consecutive terms have a constant ratio (and are nonzero).
\begin{enumerate}
\item Each term is $3$ times the previous one: \textbf{geometric}, ratio $3$. (The differences $6,18,54,\dots$ are not constant.)
\item The ratios are $2,4,16$ and the differences $2,12,240$, neither constant: \textbf{neither}. (Each term is the square of the previous one.)
\item The differences are all $4$: \textbf{arithmetic}, difference $4$.
\item Each term is twice the previous one: \textbf{geometric}, ratio $2$.
\item Each year the balance is multiplied by $1+\dfrac{3}{100}=1.03$: \textbf{geometric}, ratio $1.03$. (With simple interest, where the same amount $3\%$ of the original deposit is added every year, the balances would instead be arithmetic.)
\end{enumerate}
